\section{Discussion}
\label{sec:discussion}
ARTPS is best read as an operational architecture: anomaly cues and relative depth feed
fusion, localization, suppression, diversity, and buffering under fail-closed gates.
The design goal is structured prioritization under rover constraints, not unconstrained
claim of superior detection in all regimes.

Relative depth is an ordering cue~\citep{ranftl2021dpt}; interpreting it as metric distance
would overstate what monocular estimation provides.
Operational suppressions address common imagery distractors, but residual false positives
remain an open operational risk on the still-closed independent test split.

% TODO(results): discuss accepted-abstract reproduction only if C05--C07 close beyond pending.
% Avoid ``outperforms'' / ``demonstrate'' language until evidence gates allow.

The separation between historical accepted-abstract claims and the current independent-eval
track is intentional: conflating them would misrepresent evidence status.
The initial \texttt{independent\_eval\_v1} binary labels were produced by a classical-image
heuristic with forced $180/180$ class balance.
A model-blind repeat-author review changed the semantic distribution substantially
($340$ positive / $20$ negative).
\texttt{independent\_eval\_v1\_1} therefore follows the intended science-interest annotation
definition more faithfully on this selected domain; the imbalance indicates that
science-interest prevalence is naturally high under that definition.
Frozen-score AUROC~$0.772$ indicates useful ranking separation.
Threshold-based F1 is not informative here because the predeclared operating point is
all-positive.
A future balanced challenge set or harder negative set may be useful for operating-point
evaluation.
This label-audit finding is confined to the supplementary independent-eval benchmark; it
does not imply that the historical accepted-abstract ARTPS experiment suffered the same
issue.
Proxy analyses remain secondary and must not migrate into abstract headlines.

Safety language stays bounded: ARTPS is safety-aware software on a flight-readiness-oriented
path; it is not flight-qualified and does not assert universal safety guarantees.
Hardware timing (including planned Jetson profiles) is engineering measurement, not
certification evidence.
